\documentclass[12pt,a4paper]{amsart}

\usepackage{amsmath,amssymb,amsthm}
\usepackage{mathtools}
\usepackage[utf8]{inputenc}
\usepackage[ngerman]{babel}
\usepackage{hyperref}
\usepackage{geometry}
\geometry{margin=2.5cm}
\usepackage{booktabs}

% -------------------------------------------------------
\newtheorem{theorem}{Satz}[section]
\newtheorem{lemma}[theorem]{Lemma}
\newtheorem{korollar}[theorem]{Korollar}
\newtheorem{proposition}[theorem]{Proposition}
\theoremstyle{definition}
\newtheorem{definition}[theorem]{Definition}
\newtheorem{bemerkung}[theorem]{Bemerkung}
\newtheorem{beispiel}[theorem]{Beispiel}

% -------------------------------------------------------
\newcommand{\N}{\mathbb{N}}
\newcommand{\Z}{\mathbb{Z}}
\newcommand{\Q}{\mathbb{Q}}
\newcommand{\R}{\mathbb{R}}
\newcommand{\C}{\mathbb{C}}
\newcommand{\eq}[1]{e\!\left(#1\right)}

% -------------------------------------------------------
\begin{document}

\title{Die Hardy--Littlewood-Kreismethode:\\
       Exponentialsummen, Haupt- und Nebenbögen}

\author{Michael Fuhrmann}
\address{Specialist Mathematics Project, Build 80}
\email{specialist-maths@localhost}
\date{2026}

\begin{abstract}
Die Hardy--Littlewood-\emph{Kreismethode} ist eine der mächtigsten
Techniken der analytischen Zahlentheorie.  Sie wurde von Hardy und
Ramanujan (1918) eingeführt und von Hardy und Littlewood (1920--1928) auf
additive Probleme angewendet.  Die Grundidee ist die Integraldarstellung
der Darstellungsanzahl $r_s(n)$ (eine Zahl $n$ als Summe von $s$
$k$-ten Potenzen darzustellen) über einen Konturintegral auf dem
Einheitskreis $[0,1)$.  Das Intervall wird in \emph{Hauptbögen}
(Umgebungen rationaler Zahlen mit kleinem Nenner, wo die Exponentialsumme
groß und explizit berechenbar ist) und \emph{Nebenbögen} (wo die
Exponentialsumme klein und abschätzbar ist) zerlegt.

Wir entwickeln die Methode von Grund auf: Exponentialsummen (Weyl-Summen
und Gauß-Summen), die Farey-Zerlegung, die Struktur von Haupt- und
Nebenbögen, den Hauptterm über die \emph{singuläre Reihe} und das
\emph{singuläre Integral}, sowie Weyls Differenzierungsmethode für
Nebenbogen-Schranken.
\end{abstract}

\maketitle
\tableofcontents

% ================================================================
\section{Einleitung}
% ================================================================

Ein zentrales Problem der additiven Zahlentheorie lautet: \emph{Auf wie
viele Arten kann eine positive ganze Zahl $n$ geschrieben werden als}
\[
  n \;=\; x_1^k + x_2^k + \cdots + x_s^k
\]
mit nicht-negativen ganzen Zahlen $x_i$ (Waringsches Problem), oder als
\[
  n \;=\; p_1 + p_2 + \cdots + p_r
\]
mit Primzahlen $p_i$ (Goldbach-artige Probleme)?

Die Kreismethode transformiert solche Fragen in Integrale.  Die
\emph{Erzeugende Funktion} (Exponentialsumme)
\[
  f(\alpha) \;=\; \sum_{m=0}^{M} \eq{m^k \alpha},
  \qquad \eq{x} := e^{2\pi i x},
\]
erfüllt die Schlüsselidentität:
\begin{equation}\label{eq:key}
  r_s(n) \;=\; \int_0^1 f(\alpha)^s\, \eq{-n\alpha}\, d\alpha.
\end{equation}
Dies folgt unmittelbar aus der Orthogonalität der Charaktere
$\alpha \mapsto \eq{m\alpha}$: das Integral „greift" genau den
Koeffizienten von $\eq{n\alpha}$ in $f(\alpha)^s$ heraus.

Die Kreismethode ist die Kunst, dieses Integral durch Zerlegung
von $[0,1)$ in Gebiete zu berechnen, wo $f(\alpha)$ entweder groß
(und explizit berechenbar) oder klein (und abschätzbar) ist.

\textbf{Notation.}  $\eq{x} = e^{2\pi ix}$, $\log = \ln$, $p$
bezeichnet stets eine Primzahl, $\Lambda$ die von Mangoldt-Funktion.

% ================================================================
\section{Exponentialsummen}
\label{sec:expsummen}
% ================================================================

\subsection{Weyl-Summen}

\begin{definition}[Weyl-Summe]
  Für ein Polynom $f(x) = \alpha_k x^k + \cdots + \alpha_1 x \in \R[x]$
  und $N \geq 1$ ist die \emph{Weyl-Summe}:
  \[
    W(f; N) \;=\; \sum_{n=1}^{N} \eq{f(n)}.
  \]
  Im einfachsten Fall $f(x) = \alpha x^k$ schreiben wir $W(\alpha; k; N)$.
\end{definition}

Die triviale Schranke lautet $|W(f; N)| \leq N$.  Wenn $\alpha$ weit
von Brüchen mit kleinem Nenner entfernt ist, tritt erhebliche
Auslöschung auf.

\begin{lemma}[Weyl, 1916: Geometrische Summe]\label{lem:geom}
  Für $\alpha \in \R$ und $N \geq 1$:
  \[
    \left|\sum_{n=1}^{N} \eq{n\alpha}\right|
    \;\leq\; \min\!\left(N,\; \frac{1}{2\|\alpha\|}\right),
  \]
  wobei $\|x\| = \min_{m \in \Z} |x - m|$ der Abstand zur nächsten
  ganzen Zahl ist.
\end{lemma}

\begin{proof}
  Die Summe ist eine geometrische Reihe mit Quotient $\eq{\alpha}$.
  Falls $\eq{\alpha} = 1$ (d.h.\ $\alpha \in \Z$), ergibt sie $N$.
  Sonst:
  \[
    \left|\sum_{n=1}^N \eq{n\alpha}\right|
    = \left|\frac{\eq{\alpha} - \eq{(N+1)\alpha}}{1 - \eq{\alpha}}\right|
    \leq \frac{2}{|1 - \eq{\alpha}|}
    = \frac{1}{|\sin(\pi\alpha)|}
    \leq \frac{1}{2\|\alpha\|},
  \]
  unter Verwendung von $|\sin(\pi x)| \geq 2\|x\|$.  Das Minimum mit
  $N$ liefert die Behauptung.
\end{proof}

\subsection{Gauß-Summen}

\begin{definition}[Vollständige Gauß-Summe]
  Für ganze Zahlen $a$, $q$ mit $q \geq 1$:
  \[
    G(a, q) \;=\; \sum_{m=1}^{q} \eq{\frac{a m^k}{q}}.
  \]
\end{definition}

\begin{lemma}[Gauß-Summen-Schranke]\label{lem:gauss}
  Für $\gcd(a, q) = 1$: $|G(a, q)| \leq k \cdot q^{1/2 + \varepsilon}$.
  Für $k = 2$ (quadratische Gauß-Summen): $|G(a,q)| \leq \sqrt{2q}$.
\end{lemma}

\subsection{Exponentialsummen über Primzahlen}

\begin{definition}[Primzahl-Exponentialsumme]
  \[
    S(\alpha) \;=\; \sum_{n \leq N} \Lambda(n)\, \eq{n\alpha},
  \]
  wobei $\Lambda(n) = \log p$ für $n = p^k$ (Primzahlpotenz) und $0$ sonst.
\end{definition}

% ================================================================
\section{Die Farey-Zerlegung}
\label{sec:farey}
% ================================================================

Die Farey-Zerlegung liefert eine kanonische Partition von $[0,1)$ in
Bögen um rationale Brüche $a/q$ mit kleinem Nenner.

\begin{definition}[Farey-Folge]
  Die \emph{Farey-Folge} $\mathcal{F}_Q$ der Ordnung $Q$ ist die Menge
  aller Brüche $a/q$ mit $0 \leq a \leq q \leq Q$ und $\gcd(a,q) = 1$
  (plus $0/1$ und $1/1$), in aufsteigender Reihenfolge.
\end{definition}

\begin{lemma}[Farey-Partition]\label{lem:farey}
  Die Farey-Bögen der Ordnung $Q$ sind paarweise disjunkt und überdecken
  $[0,1)$:
  \[
    [0, 1) \;=\; \bigsqcup_{q \leq Q} \bigsqcup_{\substack{a=0\\ \gcd(a,q)=1}}^{q-1}
                 \mathfrak{M}(a, q).
  \]
  Jeder Bogen $\mathfrak{M}(a, q)$ hat Länge $\sim 1/(qQ)$.
\end{lemma}

\begin{proof}
  Dies ist ein Standardresultat in der Theorie der Farey-Folgen;
  siehe Hardy--Wright~\cite{HardyWright}, Kapitel~3.  Die Disjunktheit
  folgt aus der Mediant-Eigenschaft: Zwischen zwei aufeinanderfolgenden
  Farey-Brüchen $a/q$ und $a'/q'$ liegt der Mediant $(a+a')/(q+q')$,
  und jeder Punkt wird von genau einem Bogen beansprucht.  Die
  Überdeckung folgt daraus, dass jeder Punkt $\alpha \in [0,1)$ einen
  nächstgelegenen Farey-Bruch besitzt.
\end{proof}

% ================================================================
\section{Haupt- und Nebenbögen}
\label{sec:boegen}
% ================================================================

In der Praxis verwendet man eine vereinfachte Version mit einem festen
Parameter $Q$.

\begin{definition}[Haupt- und Nebenbögen]\label{def:boegen}
  Sei $N \geq 1$ und $Q = Q(N)$ (typisch: $Q = N^\delta$ oder
  $Q = N/(\log N)^B$).  Definiere:
  \begin{align*}
    \text{Hauptbögen:}\quad
    \mathfrak{M} &\;=\; \bigcup_{q \leq Q}\, \bigcup_{\substack{a=0\\ \gcd(a,q)=1}}^{q-1}
                         \mathfrak{M}(a, q), \\
    \mathfrak{M}(a, q) &\;=\; \left\{\alpha \in [0,1) :
      \left|\alpha - \frac{a}{q}\right| \leq \frac{Q}{N}\right\}, \\
    \text{Nebenbögen:}\quad
    \mathfrak{m} &\;=\; [0,1) \setminus \mathfrak{M}.
  \end{align*}
\end{definition}

Die Philosophie der Kreismethode ist:
\begin{enumerate}
  \item Auf \textbf{Hauptbögen} kann $f(\alpha)$ durch Gauß-Summen
    mal einem glatten Integral approximiert werden.
    Dies liefert den \emph{Hauptterm} von $r_s(n)$.
  \item Auf \textbf{Nebenbögen} ist $|f(\alpha)|$ klein (durch
    Weylsche Differenzierung oder das große Sieb).
    Dies liefert einen \emph{Fehlerterm}.
\end{enumerate}

% ================================================================
\section{Hauptbogen-Analyse}
\label{sec:hauptbogen}
% ================================================================

\begin{lemma}[Hauptbogen-Approximation]\label{lem:hauptapprox}
  Für $\alpha = a/q + \beta$ mit $\gcd(a,q)=1$, $q \leq Q$ und
  $|\beta| \leq Q/N$:
  \[
    f(\alpha) \;=\; \frac{G(a, q)}{q}\, v(\beta) \;+\; O\!\left(q^{1+\varepsilon}\right),
  \]
  wobei $v(\beta) = \int_0^M \eq{\beta t^k}\, dt$ das \emph{Weyl-Integral} ist.
\end{lemma}

\subsection{Die singuläre Reihe}

\begin{definition}[Singuläre Reihe]\label{def:sing_reihe}
  \[
    \mathfrak{S}(n) \;=\; \sum_{q=1}^{\infty} \sum_{\substack{a=1\\ \gcd(a,q)=1}}^{q}
                          \left(\frac{G(a,q)}{q}\right)^s \eq{-\frac{an}{q}}.
  \]
  Diese hat die Euler-Produktdarstellung
  $\mathfrak{S}(n) = \prod_p \sigma_p(n)$,
  wobei $\sigma_p(n)$ die $p$-adische Dichte von Lösungen ist.
\end{definition}

\begin{bemerkung}
  $\mathfrak{S}(n) > 0$ genau dann, wenn die Gleichung
  $x_1^k + \cdots + x_s^k = n$ für alle Primzahlen $p$ $p$-adische
  Lösungen hat (lokale Lösbarkeit).
\end{bemerkung}

\subsection{Das singuläre Integral}

\begin{definition}[Singuläres Integral]
  \[
    \mathfrak{J}(n) \;=\; \int_{-\infty}^{\infty} v(\beta)^s\, \eq{-n\beta}\, d\beta.
  \]
\end{definition}

\begin{lemma}[Singuläres Integral, Abschätzung]\label{lem:sing_int}
  Für $s > k$ und $n > 0$:
  $\mathfrak{J}(n) \sim C_{k,s} \cdot n^{s/k - 1}$
  mit einer von $k$ und $s$ abhängigen Konstante $C_{k,s} > 0$.
\end{lemma}

\begin{theorem}[Hauptbogen-Beitrag]\label{thm:hauptbogen}
  Für $s > 2k+1$:
  \[
    \int_{\mathfrak{M}} f(\alpha)^s \eq{-n\alpha}\, d\alpha
    \;=\; \mathfrak{S}(n)\, \mathfrak{J}(n)\, N^{s/k - 1}
          \;+\; O\!\left(N^{s/k - 1 - \delta}\right)
  \]
  für ein $\delta = \delta(k, s) > 0$.
\end{theorem}

% ================================================================
\section{Nebenbogen-Schranken: Weyls Ungleichung}
\label{sec:nebenbogen}
% ================================================================

\begin{theorem}[Weylsche Ungleichung]\label{thm:weyl}
  Sei $k \geq 2$ und $|\alpha - a/q| \leq 1/q^2$ für ganze Zahlen
  $a, q$ mit $\gcd(a,q)=1$ und $1 \leq q \leq N^k$.  Dann gilt:
  \[
    \left|\sum_{n=1}^{N} \eq{n^k \alpha}\right|
    \;\ll\; N^{1+\varepsilon}
             \left(\frac{1}{q} + \frac{1}{N} + \frac{q}{N^k}\right)^{2^{1-k}}.
  \]
\end{theorem}

\begin{proof}[Beweisidee (Weylsche Differenzierung für $k=2$)]
  Wir quadrieren die Summe:
  \[
    \left|\sum_n \eq{n^2 \alpha}\right|^2
    = \sum_h \sum_n \eq{h(2n + h)\alpha},
  \]
  wobei $h = n - m$.  Die innere Summe über $n$ ist eine geometrische
  Reihe, beschränkt durch $\min(N, \|2h\alpha\|^{-1})$.  Summation
  über $h \leq N$ liefert die Verbesserung gegenüber der trivialen Schranke.
  Für allgemeines $k$ iteriert man die Differenzierung $k-1$-mal.
\end{proof}

\begin{korollar}[Nebenbogen-Schranke]\label{cor:nebenbogen}
  Auf den Nebenbögen $\mathfrak{m}$ gilt:
  \[
    |f(\alpha)| \;\ll\; N^{1 - \sigma(k) + \varepsilon},
  \]
  wobei $\sigma(k) = 2^{1-k}$ (Weyl) oder besser (Vinogradov--Wooley).
\end{korollar}

% ================================================================
\section{Die vollständige Kreismethode}
\label{sec:vollstaendig}
% ================================================================

\begin{theorem}[Hardy--Littlewood--Vinogradov]\label{thm:haupt}
  Für $k \geq 2$ und $s \geq s_0(k)$ und alle hinreichend großen $n$:
  \[
    r_s(n) \;=\; \mathfrak{S}(n)\, \mathfrak{J}(n)\, n^{s/k - 1}
                 \;+\; O\!\left(n^{s/k - 1 - \delta}\right).
  \]
  Insbesondere gilt $r_s(n) > 0$ für alle hinreichend großen $n$.
\end{theorem}

\begin{proof}[Beweis (Skizze)]
  Man zerlegt:
  $r_s(n) = \int_{\mathfrak{M}} f^s \eq{-n\cdot} + \int_{\mathfrak{m}} f^s \eq{-n\cdot}$.

  \textbf{Hauptbogen:} Satz~\ref{thm:hauptbogen} liefert den Hauptterm.

  \textbf{Nebenbogen:} Mit Korollar~\ref{cor:nebenbogen} und Parsevals Identität:
  \[
    \left|\int_{\mathfrak{m}} f^s \eq{-n\cdot}\right|
    \leq N^{(s-2)(1-\sigma)} \cdot N \;=\; N^{s-1-(s-2)\sigma}.
  \]
  Für $s > 1 + 2/\sigma$ ist dies $o(N^{s/k - 1})$.
\end{proof}

% ================================================================
\section{Anwendung auf Goldbach-artige Probleme}
\label{sec:goldbach}
% ================================================================

Für Primzahl-Summen verwendet man
$S(\alpha) = \sum_{n \leq N} \Lambda(n) \eq{n\alpha}$.

\begin{lemma}[Hauptbogen für Primzahlen]\label{lem:hauptbogen_primes}
  Für $\alpha = a/q + \beta$ mit kleinem $q$:
  $S(\alpha) \approx \frac{\mu(q)}{\phi(q)}\, v_N(\beta)$.
\end{lemma}

Der Faktor $\mu(q)/\phi(q)$ kodiert den Beitrag des Hauptcharakters mod $q$.
Die Nebenbogen-Schranke für $S(\alpha)$ folgt aus dem Bombieri--Vinogradov-Satz
und dem großen Sieb (vgl.\ Paper~15):
\[
  \max_{\alpha \in \mathfrak{m}} |S(\alpha)| \;\ll\; N(\log N)^{-B}.
\]

% ================================================================
\section{Historische Anmerkungen und Erweiterungen}
\label{sec:hist}
% ================================================================

\begin{bemerkung}[Geschichte]
  Die Kreismethode wurde erstmals von Hardy und Ramanujan (1918) zur
  asymptotischen Formel für die Partitionsfunktion $p(n)$ verwendet.
  Hardy und Littlewood (1920--1923) wendeten sie auf das Waringsche
  Problem an.  Vinogradov (1937) verbesserte die Nebenbogen-Schranken
  drastisch.  Wooley (1992--2016) bewies die Vinogradov-Hauptvermutung
  durch \emph{effizientes Kongruenzrechnen} und erhielt die bisher
  besten Schranken für $G(k)$.
\end{bemerkung}

\begin{bemerkung}[Grenzen der Kreismethode]
  Die Kreismethode hat fundamentale Grenzen:
  \begin{itemize}
    \item Sie benötigt $s \geq s_0(k)$ „genug" Summanden.  Für binäres
      Goldbach ($r = 2$ Primzahlen) liefert sie nur bedingte Resultate.
    \item Das \emph{Paritätsproblem} der Siebtheorie erklärt, warum
      $s = 2$ außer Reichweite liegt.
    \item Für $r \geq 3$ Primzahlen gelingt der Beweis unbedingt
      (Vinogradovs Satz, Paper~18).
  \end{itemize}
\end{bemerkung}

% ================================================================
\section{Schlussfolgerung}
% ================================================================

Die Hardy--Littlewood-Kreismethode beruht auf einer einzigen eleganten
Idee: der Integraldarstellung~\eqref{eq:key} der Darstellungsanzahl.

\begin{enumerate}
  \item In der Nähe eines Bruchs $a/q$ mit kleinem $q$
    (\textbf{Hauptbögen}) wird die Exponentialsumme durch
    Gauß-Summen mal einem Integral approximiert.  Die Summe über
    alle solchen Brüche liefert die \emph{singuläre Reihe} $\mathfrak{S}(n)$
    und das \emph{singuläre Integral} $\mathfrak{J}(n)$.
  \item Fernab von Brüchen mit kleinem Nenner (\textbf{Nebenbögen})
    ist die Exponentialsumme nach Weyls Ungleichung klein.
    Das Nebenbogen-Integral ist ein Fehlerterm.
\end{enumerate}

Die folgenden Papers (18--20) wenden und erweitern diese Methode.

% ================================================================
\begin{thebibliography}{9}

\bibitem{HardyRamanujan1918}
G.~H.~Hardy und S.~Ramanujan.
\newblock Asymptotic formulae in combinatory analysis.
\newblock \emph{Proc.\ London Math.\ Soc.\ (2)}, 17:75--115, 1918.

\bibitem{HardyLittlewood1923}
G.~H.~Hardy und J.~E.~Littlewood.
\newblock Some problems of ``Partitio Numerorum''. III.
\newblock \emph{Acta Math.}, 44:1--70, 1923.

\bibitem{Vaughan1981}
R.~C.~Vaughan.
\newblock \emph{The Hardy--Littlewood Method}, 2.~Aufl.
\newblock Cambridge University Press, 1997.

\bibitem{Vinogradov1954}
I.~M.~Winogradow.
\newblock \emph{Die Methode der Trigonometrischen Summen in der Zahlentheorie}.
\newblock Harri Deutsch, Frankfurt, 1989. (Deutsche Übersetzung.)

\bibitem{Wooley2016}
T.~D.~Wooley.
\newblock The cubic case of the main conjecture in Vinogradov's mean value theorem.
\newblock \emph{Adv.\ Math.}, 294:532--561, 2016.

\bibitem{IwaniecKowalski}
H.~Iwaniec und E.~Kowalski.
\newblock \emph{Analytic Number Theory}.
\newblock American Mathematical Society, Providence, RI, 2004.

\end{thebibliography}

\end{document}
